% =========================================================
% Capstone Exercise — Lambda Calculus
% =========================================================

\subsection*{Capstone: The Identity Proof Under Curry--Howard}
\label{prf:capstone-lambda-calculus}

\begin{remark*}[Return]
Chapter capstone --- synthesises the introduction section.
\end{remark*}

\begin{proof}
The claim is threefold:
\begin{enumerate}
  \item The lambda term $I := \lambda x.x$ is well-typed: in the
        empty context, $\vdash \lambda x.x : A \to A$ for any type $A$.
  \item Under the Curry--Howard correspondence, this typing derivation
        \emph{is} the proof of the tautology $A \Rightarrow A$.
  \item In Lean~4, the term \texttt{fun x => x} proves any goal
        of the form \texttt{A \to\ A}, and \texttt{\#check @id} confirms
        the type is \texttt{$\forall$ (A : Type), A \to\ A}.
\end{enumerate}

% TODO: Work out parts (1) and (2) by hand.
%
% Part (1): Construct the typing derivation tree using the three
% typing rules from Definition~\ref{def:typing-judgment}.
% The derivation has exactly two steps: (Var) then (Abs).
%
% Part (2): Identify each rule in the derivation with its logical
% counterpart: (Var) = hypothesis; (Abs) = implication introduction.
%
% Part (3): Write the Lean term and verify it compiles.
% This is a paper exercise; no Lean file is generated.
\end{proof}

\begin{remark*}[Proof shape]
% TODO: fill once worked by hand.

\FlashProof{1. Apply (Var): $x:A \vdash x:A$.
  2. Apply (Abs): $\vdash \lambda x.x : A\to A$.
  3. Identify with logical proof: hypothesis then implication introduction.}
\end{remark*}

\begin{remark*}[Dependencies]
\begin{itemize}
  \item \hyperref[def:typing-judgment]{Definition of typing judgment}:
        the three typing rules
  \item \hyperref[thm:curry-howard]{Curry--Howard correspondence}:
        the dictionary between types and propositions
\end{itemize}

\FlashUses{def:typing-judgment, thm:curry-howard}
\end{remark*}
